\documentclass[a4paper,landscape,headrule,footrule]{foils}
%%
%%% macros for Theories of Grammar
%%%
\usepackage{polyglossia}
\setdefaultlanguage{english}
%\setmainfont{TeX Gyre Pagella}


\newcommand{\logo}{~}
\newcommand{\header}[3]{%
\title{\vspace*{-2ex} \large JMORF/MORS --- Morpho-Syntax
\\[2ex] \Large  \red{#2} \\ \red{#3}}
\author{\blu{Francis Bond}   \\ 
\normalsize  \textbf{Palacký University}\\
\normalsize  \url{https://fcbond.github.io/}\\
\normalsize  \texttt{bond@ieee.org}}
\MyLogo{JMORF (2026)}
\renewcommand{\logo}{#2}
\hypersetup{
   pdfinfo={
     Author={Francis Bond},
     Title={#1: #2},
     Subject={JMORF/MORS: Morpho-Syntax},
     Keywords={Syntax, Morphology, HPSG, Unification, Constructions},
     License={CC BY 4.0}
   }
 }

\date{#1
  \\ Location: KB-1.11}
}
\usepackage{hyperref}
\hypersetup{
     colorlinks,
     linkcolor={blue!50!black},
     citecolor={red!50!black},
     urlcolor={blue!80!black}
   }



\usepackage{xcolor}
\usepackage{graphicx}
\newcommand{\blu}[1]{\textcolor{blue}{#1}}
\newcommand{\red}[1]{\textcolor{red}{#1}}
\newcommand{\grn}[1]{\textcolor{green}{#1}}
\newcommand{\hide}[1]{\textcolor{white}{#1}}
\newcommand{\emp}[1]{\textcolor{blue}{#1}}
\newcommand{\txx}[1]{\textbf{\textcolor{blue}{#1}}}
\newcommand{\lex}[1]{\textbf{\mtcitestyle{#1}}}


% \renewcommand{\labelitemi}{\textcolor{violet}{‣}}
% \renewcommand{\labelitemii}{\textcolor{purple}{‣}}

\newcommand{\subhead}[1]{\noindent\textbf{#1}\\[5mm]}

\newcommand{\Bad}{\emp{\raisebox{0.15ex}{\ensuremath{\mathbf{\otimes}}}}}
\newcommand{\bad}[1]{*\eng{#1}}

\newcommand{\com}[1]{\hfill (\emp{#1})}%

\usepackage{relsize,xspace}
\newcommand{\into}{\ensuremath{\rightarrow}\xspace}
\newcommand{\tot}{\ensuremath{\leftrightarrow}\xspace}
\usepackage{url}
\newcommand{\lurl}[1]{\MyLogo{\url{#1}}}

\usepackage{langsci-gb4e}
\let\eachwordone=\slshape
%%% langsci-gb4e defines \thexnumi as \@roman\@xsi{xnumi}, which only works
%%% inside an exe environment; some decks use xlisti on its own.  Restore
%%% the plain gb4e definition so that keeps working.
\makeatletter
\def\thexnumi{\@xsi{xnumi}}
\makeatother
\newcommand{\lx}[1]{\textbf{\mtciteform{#1}}}
\newcommand{\ix}{\ex\slshape}



\newcommand{\ent}{\ensuremath{\Rightarrow}\xspace}
\newcommand{\ngv}{\ensuremath{\not\Rightarrow}\xspace}
%\usepackage{times}
%\usepackage{nttfoilhead}
\newcommand{\myslide}[1]{\foilhead[-25mm]{\raisebox{12mm}[0mm]{\red{#1}}}\MyLogo{\logo}}
\newcommand{\myslider}[1]{\rotatefoilhead[-25mm]{\raisebox{12mm}[0mm]{\red{#1}}}}
%\newcommand{\myslider}[1]{\rotatefoilhead{\raisebox{-8mm}{\emp{#1}}}}

\newcommand{\section}[1]{\myslide{}{\begin{center}\Huge \red{#1}\end{center}}}



\usepackage[lyons,j,e,k]{mtg2e}
%\renewcommand{\mtcitestyle}[1]{\textcolor{-red!75!green!50}{\textsl{#1}}}
\renewcommand{\mtcitestyle}[1]{\textcolor{teal}{\textsl{#1}}}
\newcommand{\cs}{\mtciteform}
\newcommand{\iz}[1]{\texttt{\textup{#1}}}
\newcommand{\gm}{\textsc}
\usepackage[normalem]{ulem}
\newcommand{\ul}{\uline}
\newcommand{\ull}{\uuline}
\newcommand{\wl}{\uwave}
\newcommand{\vs}{\ensuremath{\Leftrightarrow}~}
%%%
%%% Bibliography
%%%
\usepackage{natbib}
%\usepackage{url}
\usepackage{bibentry}



\usepackage[utf8]{inputenc}
%%% trees
\usepackage{multicol}
%%% pst-node: one example in lec-12-control draws a connector between
%%% two points with \rnode/\ncbar.  This used to come in with rtrees.
\usepackage{pst-node}
\usepackage{forest}
\useforestlibrary{linguistics}
\forestapplylibrarydefaults{linguistics}

%%% \makebox's optional arguments are read as AVM delimiters by
%%% langsci-avm, so hide them behind a macro with no visible brackets.
\newcommand{\lbox}[1]{\makebox[0mm][l]{#1}}

%%% Tree styles.
%%%   words  -- trees whose leaves are actual words (italicised, aligned)
%%%   cat    -- trees whose leaves are categories or feature structures
\forestset{
  words/.style={for tree={s sep=7mm, l sep=5mm,
                          if n children=0{font=\itshape, tier=terminal}{}}},
  cat/.style={for tree={s sep=7mm, l sep=5mm}},
}
\newcommand{\wtree}[1]{\begin{forest} words #1 \end{forest}}
\newcommand{\ctree}[1]{\begin{forest} cat #1 \end{forest}}
%%% the rule being applied at this step of a derivation
\newcommand{\step}[1]{\begin{center}\emp{#1}\end{center}}

%%% AVMs (Language Science Press)
\usepackage{langsci-avm}

%%% Macros the decks use, carried over from headx.tex
\newcommand{\blank}{\rule{3em}{1pt}\xspace}
\newcommand{\ft}[1]{\textsc{#1}}
\newcommand{\idn}{\mtciteform}    %%% Indonesian
\newcommand{\prd}[1]{\textbf{#1}}
\newcommand{\typ}[1]{\textit{#1}}
%%% \val: provide before renewing, in case a package defines it.
\providecommand{\val}[1]{}
\renewcommand{\val}[1]{\textit{#1}}
%%% Manning's \avmbox, still used in tree labels and in maths where a
%%% bare langsci \tag would not work.  \avm{\tag{...}} gives the same
%%% boxed tag as the ones inside feature structures.
\newcommand{\avmbox}[1]{\mbox{\avm{\tag{#1}}}}

%%% Ring a value in a dotted oval.  These used pstricks \ovalnode,
%%% which came in with rtrees; tikz (via forest) does the same job.
\newcommand{\OV}[1]{\tikz[baseline]{\node[draw,dotted,red,ellipse,
          inner sep=1pt,outer sep=0pt]{#1};}}
\newcommand{\OVB}[1]{\tikz[baseline]{\node[draw,dotted,blue,ellipse,
          inner sep=1pt,outer sep=0pt]{#1};}}

%%% From CSLI book
\newcommand{\mc}{\multicolumn}
\newcommand{\HD}{\textbf{H}\xspace}
\newcommand{\el}{< >}       %%% empty list; langsci-avm needs the space

\newcommand{\A}{\noindent\textbf{A}: }
\newcommand{\Q}{\noindent\textbf{Q}: }

%%% Question and answer on two pages, written once:
%%%   \qaslide{Title}{... question ... \answer{... answer ...}}
%%% prints the slide twice, first without the \answer parts and then
%%% with them, so the answer appears "on click".  The example counter
%%% is reset so both pages show the same example numbers: gb4e keeps the
%%% running number in the equation counter between exe environments, and
%%% in xnumi inside one.
\newif\ifshowanswer
\newcommand{\answer}[1]{\ifshowanswer#1\fi}
\newcounter{qaxnumi}
\newcounter{qaeqn}
\newcommand{\qaslide}[2]{%
  \setcounter{qaxnumi}{\value{xnumi}}%
  \setcounter{qaeqn}{\value{equation}}%
  \showanswerfalse\myslide{#1}#2%
  \setcounter{xnumi}{\value{qaxnumi}}%
  \setcounter{equation}{\value{qaeqn}}%
  \showanswertrue\myslide{#1}#2%
  \showanswerfalse}
%\newcommand{\C}{\noindent\textbf{C}: }



%\usepackage{times}
%\usepackage{bchart}
%\usepackage{pgfplots}
\begin{document}
\header{Lecture 10}{Raising and Control}{}
\maketitle

%


\myslide{Overview}
\MyLogo{\citet{sag:wasow:bender:2003} --- Chapter 12}
\begin{itemize}
\item Introduction: two kinds of NP-V-(NP)-\lex{to}-VP
\item Infinitival \lex{to}
\item Subject-raising verbs
\item Subject-control verbs
\item Raising and control in Transformational Grammar
\item Object raising and object control
\item Raising and control in Czech
\item Problem Set 10: SWB 12.1, 12.2, 12.4, 12.6
%\item If time: Problem 12.4
\end{itemize}

\section{Introduction}

\qaslide{Where We Are \& Where We’re Going}{

\begin{itemize}
\item Revision
  \begin{itemize}
  \item In the last two lectures, we have seen a kind of 
    \txx{subject sharing} --- that is, cases where one NP 
    served as the \ft{spr} for two different verbs.  
    Examples?
  \answer{\\ {\small\emp{Copula and passive \lex{be} (\eng{Pat is tall}, \eng{Kim was arrested}), existential \lex{be}.}}}
  \item Last time, we looked at \txx{dummy} NPs --- that is, 
    non-referential NPs.  Examples?
  \answer{\\ {\small\emp{\lex{there}, \lex{it} (\eng{It is raining}), \lex{advantage}, \lex{tabs}.}}}
  \end{itemize}
\item Today, we’re going to look at the kind of subject 
sharing we saw with \lex{be} in more detail.
\item Then we’ll look at another kind of subject 
sharing, using dummy NPs in differentiating the 
two kinds.
\end{itemize}
}

\myslide{What Makes This Topic Different}
\begin{itemize}
\item The phenomena we have looked at so far (\txx{agreement},
  \txx{binding}, \txx{imperatives}, \txx{passives},
  \txx{existentials}, \txx{extraposition}) are easy to pick out on the
  basis of their form alone.
\item In this chapter, we look at constructions with the general form
  NP-V-(NP)-\lex{to}-VP.  It turns out that they divide into two
  kinds, differing in both syntactic and semantic properties.
\end{itemize}

\myslide{The Central Idea}

\begin{itemize}
\item  \eng{Pat continues to avoid conflict} and 
\eng{Pat tries to avoid conflict} 
both have the form NP-V-\lex{to}-VP
  \begin{itemize}
  \item But \lex{continue} is semantically a one-place 
    predicate, expressing a property of a situation 
    (namely, that it continues to be the case)
    \\ \prd{continue(avoid(Pat,conflict))}
  \item Whereas \lex{try} is semantically a two-place 
    predicate, expressing a relation between someone 
    who tries and a situation s/he tries to bring about.
    \\ \prd{try(Pat, avoid(Pat, conflict))}
  \end{itemize}
\item This semantic difference has syntactic effects.
\end{itemize}

\section{Infinitival \lex{to}}

\myslide{The Status of Infinitival \lex{to}}


\begin{itemize}
\item It’s not obvious what part of speech to assign to \lex{to}.  
\item It’s not the same as the preposition \lex{to}:
  \begin{exe}
    \ex \eng{Pat aspires to stardom.}
    \ex \eng{Pat aspires to be a good actor.}
    \ex *?\eng{Pat aspires to stardom and to be a good actor.}
    \ex \bad{Pat aspires to stardom and be a good actor.}
  \end{exe}
\item We call it an \txx{auxiliary verb}, because this will make 
our analysis of auxiliaries a little simpler.
\end{itemize}

\myslide{The Lexical Entry for Infinitival \lex{to}}
\begin{center}
  \begin{small}
    \avm{< \textnormal{to}, [
      syn &  [ head & [ form & base\\
      inf & $+$ \\
      aux & $+$ ] ] \\
      arg-st & < \1,
      [ head & [{\tiny \it verb} \\
      inf & $-$ \\
      form & base ] \\
      val & [ spr & < \1 > \\
      comps & < ~ > ] \\
      sem & [ index & s ] ] > \\
      sem & [ index & s \\ restr & < > ] ] >
}
  \end{small}
\end{center}
\myslide{The Syntax of Infinitival  \lex{to}}

\begin{center}
  \avm{[ syn &  [ head & [ form & base\\
    inf & $+$ \\
    aux & $+$ ] ] ]
}
\end{center}

\begin{itemize}\addtolength{\itemsep}{-1ex}
\item This makes it a verb, because \ft{aux} is declared on \val{verb}
\item {[\ft{inf}  $+$]} uniquely identifies the infinitival \lex{to}
\item Verbs select complements with different combinations 
  of \ft{form} and \ft{inf} values, e.g.
  \begin{itemize}
  \item complements of \lex{condescend} are [\ft{form} \val{base}] and [\ft{inf} $+$]
  \item complements of \lex{should} are [\ft{form} \val{base}] and [\ft{inf} $-$]
  \item complements of \lex{help} are [\ft{form} \val{base}]
  \end{itemize}
\item The meaning of [\ft{aux} $+$] becomes clear in Chapter 13
\end{itemize}

\qaslide{The Argument Structure}{
\begin{center}
  \avm{[ arg-st & < \1,
      [ head & [{\tiny \it verb} \\
      inf & $-$ \\
      form & base ] \\
      val & [ spr & < \1 > \\
      comps & < ~ > ] \\
      sem & [ index & s ] ] > ]
}
\end{center}
\begin{itemize}
\item What kind of constituent is the second argument?
  \answer{\\ {\small\emp{A base-form VP, not itself headed by \lex{to} ([\ft{inf} $-$]): \eng{(to) \ul{be a good actor}}.}}}
\item The tagging of the first argument and the \ft{spr} of the second 
argument is exactly like \lex{be}.
\end{itemize}
}

\qaslide{The Semantics of Infinitival \lex{to}}{
\begin{center} \small
  \avm{[ arg-st & < \1,
      [ head & [{\tiny \it verb} \\
      inf & $-$ \\
      form & base ] \\
      val & [ spr & < \1 > \\
      comps & < ~ > ] \\
      sem & [ index & \OV{s} ] ] > \\
      sem & [ index & \OV{s} \\ restr & < > ] ]
}
\end{center}

\begin{itemize}
\item The \ft{index} value is taken from the \ft{sem} of the second
  argument.
\item What is the semantic contribution of \lex{to}?
  \answer{\\ {\small\emp{None: its \ft{restr} is empty, and its \ft{index} is its complement's.}}}
\end{itemize}
}

\section{Subject-Raising Verbs}

\myslide{Dummies and \lex{continue}}

\begin{itemize}
\item Some examples:
  \begin{exe}
\ex \eng{There continue to be seats available.}
\ex \eng{It continues to matter that we lost.}
\ex \eng{Advantage continues to be taken of the innocent.}
\ex \bad{It continues to be seats available.}
\ex \bad{There continues to matter that we lost.}
\ex \bad{Advantage continues to be kept of the innocent.}
\end{exe}
\item Generalization:  Non-referential NPs can appear as the 
subject of \lex{continue} just in case they could be the subject 
of the complement of \lex{continue}.
\end{itemize}

\myslide{A New Type, for Verbs like \lex{continue}}

\begin{center} \small
\txx{Subject-Raising Verb Lexeme} (\val{srv-lxm}) \\
  \avm{[ arg-st & < \1,
      [ 
      val & [ spr & < \1 > \\
      comps & < ~ > ] \\
      sem & [ index & \OV{s} ] ] > \\
      sem & [ restr & < [ arg & s ] > ] ]
}
\end{center}

\begin{itemize}\addtolength{\itemsep}{-1ex}
\item The subject sharing is just like for \lex{be} and \lex{to}: the
  subject of \lex{continue} is also the subject of its complement
\item  \lex{continue} imposes no other constraints on its subject
\item The index of the complement must be an argument of the 
predication introduced by the verb
\end{itemize}

\myslide{The Lexical Entry for \lex{continue}}

\bigskip
\begin{center}
    \avm{< \textnormal{continue}, [{\tiny \val{srv-lxm}}\\
      arg-st & < X,
                  VP [ inf & $+$ ] > \\
      sem & [ index & s$_1$ \\ 
               restr & < [ reln & \textbf{continue} \\ 
                             sit & s$_1$ ] > ] ] >
}
\end{center}



\myslide{\lex{continue} with Inherited Information}

\begin{center}
  \begin{small}
    \avm{< \textnormal{continue}, [{\tiny \val{srv-lxm}}\\
      syn &  [ head & [{\tiny \val{verb}} \\
      pred & $-$ \\
      inf & $-$ \\
      agr & \2 ] ] \\
      arg-st & < \1 [ head & \val{nominal}\\
                         val & [ spr & < > \\ comps & < > ]  ],
                  VP [ inf & $+$ \\ spr  & < \1 > \\ index & s$_2$ ]
      > \\
      sem & [ mode & prop \\ index & s$_1$ \\ 
               restr & < [ reln & \textbf{continue} \\ 
                             sit & s$_1$ \\ arg & s$_2$  ] > ] ] >
}
  \end{small}
\end{center}


\myslide{Key Property of Subject-Raising Verbs}

The subject plays no semantic role in the predication introduced by
the \txx{SRV} itself.  Its semantic role (if any) is only in the
predication introduced in the complement.

\begin{center}
  \begin{small}
    \avm{< \textnormal{continue}, [{\tiny \val{srv-lxm}}\\
      arg-st & < \OV{\1} [ head & \val{nominal}\\
                         val & [ spr & < > \\ comps & < > ]  ],
                  VP [ inf & $+$ \\ spr  & < \OV{\1} > \\ 
                        index & \OVB{s$_2$} ]
      > \\
      sem & [ mode & prop \\ index & s$_1$ \\ 
               restr & < [ reln & \textbf{continue} \\ 
                             sit & s$_1$ \\ arg & \OVB{s$_2$}  ] > ] ] >
}
  \end{small}
\end{center}


\myslide{Constraints on SRVs' Subjects Are from Their Complements}

\begin{itemize}
\item SRVs take dummy subjects when and only when their 
complements do.
\item SRVs take idiom chunk subjects when and only when their 
complements do.
\item Passivizing the complement of an SRV doesn’t change the 
truth conditions of the whole sentence:
\begin{exe}
\ex \eng{Skeptics continue to question your hypothesis.}
\ex \eng{Your hypothesis continues to be questioned by skeptics.}
\end{exe}
\end{itemize}

\myslide{\lex{continue} with Active Complement}
\begin{minipage}[c]{0.58\textwidth}
\scalebox{0.72}{\begin{forest}
cat
[{S\avmbox{A}} [{\avmbox{1} NP$_i$} [{Skeptics}]]
[{VP \avm{[ \ft{spr} < \avmbox{1} > ]}} [{V \avm{[ \ft{spr} < \avmbox{1} > ]}} [{continue}]]
[{VP \avm{[ \ft{spr} < \avmbox{1} > ]}} [{V \avm{[ \ft{spr} < \avmbox{1} > ]}} [{to}]]
[{VP \avm{[ \ft{spr} < \avmbox{1} > ]}} [{V \avm{[ \ft{spr} < \red{\avmbox{1}} > ]}} [{question}]]
[{NP$_j$} [{your hypothesis}, roof]]]]]]
\end{forest}}
\end{minipage}\hfill
\begin{minipage}[c]{0.4\textwidth}
\scalebox{0.8}{\avm{[ restr & \avmbox{A} < [ reln & \textbf{continue} \\  sit & s$_1$ \\ arg & s$_2$ ], \\
        [ reln & \textbf{question} \\  sit  & s$_2$ \\ doubter & \OV{i} \\ doubted & j],
  \ldots > ]}}
\end{minipage}

\emp{The subject \avmbox{1} is passed down unchanged to \lex{question},
  whose \ft{spr} gives it the \ft{doubter} role.  \lex{continue} itself
  gives it no role.}

\myslide{\lex{continue} with Passive Complement}
\begin{minipage}[c]{0.58\textwidth}
\scalebox{0.6}{\begin{forest}
cat
[{S\avmbox{A}} [{\avmbox{1} NP$_j$} [{Your hypothesis}, roof]]
[{VP \avm{[ \ft{spr} < \avmbox{1} > ]}} [{V \avm{[ \ft{spr} < \avmbox{1} > ]}} [{continues}]]
[{VP \avm{[ \ft{spr} < \avmbox{1} > ]}} [{V \avm{[ \ft{spr} < \avmbox{1} > ]}} [{to}]]
[{VP \avm{[ \ft{spr} < \avmbox{1} > ]}} [{V \avm{[ \ft{spr} < \avmbox{1} > ]}} [{be}]]
[{VP \avm{[ \ft{spr} < \avmbox{1} > ]}} [{V \avm{[ \ft{spr} < \red{\avmbox{1}} > ]}} [{questioned}]]
[{PP$_i$} [{P$_i$} [{by}]]
[{NP$_i$} [{skeptics}]]]]]]]]
\end{forest}}
\end{minipage}\hfill
\begin{minipage}[c]{0.4\textwidth}
\scalebox{0.8}{\avm{[ restr & \avmbox{A} <[ reln & \textbf{continue} \\  sit & s$_1$ \\ arg & s$_2$ ], \\
        [ reln & \textbf{question} \\ sit  & s$_2$ \\ doubter & i \\ doubted & \OV{j}],
  \ldots > ]}}
\end{minipage}

\emp{Now the passive participle gives the subject the \ft{doubted}
  role.  Same \textbf{question} predication, same truth conditions:
  skeptics doubt, the hypothesis is doubted.}

\section{Subject-Control Verbs}

\myslide{Control Verbs}

\begin{itemize}
\item Control verbs, like \lex{try}, appear in contexts that 
look just like the contexts for raising verbs:
\begin{xlisti}
\ex \eng{Pat tried to stay calm.}
\trans looks superficially like
\ex \eng{Pat continued to stay calm.}
\end{xlisti}
\item Control verbs also share their subjects with their 
complements, but in a different way.
\item A control verb expresses a relation between the 
referent of its subject and the situation denoted by 
its complement.
\end{itemize}

\myslide{Control Verbs Are Not Transparent }


\begin{itemize}
\item They never take dummies or idiom chunks as 
  subjects.
  \begin{xlisti}
    \ex \bad{There try to be bugs in my program.}
    \ex \bad{It tries to upset me that the Giants lost.}
    \ex \bad{Advantage tries to be taken of tourists.}
  \end{xlisti}
\item Passivizing the complement’s verb changes the truth 
conditions.
\begin{xlisti}
    \ex \eng{The police tried to arrest disruptive demonstrators.}
    \trans $\not=$
    \ex \eng{Disruptive demonstrators tried to be arrested by the police.}
  \end{xlisti}
\end{itemize}

\myslide{A New Type}
\begin{center} \small
\txx{Subject-Control Verb Lexeme} (\val{scv-lxm})\\
  \avm{[ arg-st & < NP$_i$,
      [ 
      val & [ spr & < NP$_i$ > \\
      comps & < ~ > ] \\
      sem & [ index & \OV{s} ] ] > \\
      sem & [ restr & < [ arg & s ] > ] ]
}
\end{center}

\begin{itemize}\addtolength{\itemsep}{-2ex}
\item This differs from \val{srv-lxm} in that the first argument and the
   \ft{spr} of the second argument are coindexed, not tagged (re-entrant)
\item This means that they only need to share \ft{index} values, but may
   differ on other features
\item And the first argument --- the subject --- must have an \ft{index}
  value, so it cannot be non-referential
\end{itemize}

\myslide{The Lexical Entry for \lex{try}}

\bigskip
\begin{center}
    \avm{< \textnormal{try}, [{\tiny \val{scv-lxm}}\\
      arg-st & < NP$_i$,
                  VP [ inf & $+$ ] > \\
      sem & [ index & s$_1$ \\ 
               restr & < [ reln & \textbf{try} \\ 
                             sit & s$_1$ \\
                             trier & i ] > ] ] >
}
\end{center}


Note that the subject (NP$_i$) plays a semantic role with 
respect to the verb, namely the \ft{trier}.

\myslide{\lex{try} with Inherited Information}

\begin{tabular}[t]{cc}
\begin{minipage}{0.5\linewidth}
  \footnotesize
    \avm{< \textnormal{try}, [{\tiny \val{scv-lxm}}\\
      syn &  [ head & [{\tiny \val{verb}} \\
      pred & $-$ \\
      inf & $-$ \\
      agr & \1 ] ] \\
      arg-st & < NP$_i$,
                  VP [ inf & $+$ \\ spr  & < NP$_i$ > \\ index & s$_2$ ]
      > \\
      sem & [ mode & prop \\ index & s$_1$ \\ 
               restr & < [ reln & \textbf{try} \\ 
                             sit & s$_1$ \\ 
                             arg & s$_2$\\
                             trier & i  ] > ] ] >
}
  \end{minipage}
  &
\begin{minipage}{0.4\linewidth}
\small
Things to Note:
\begin{itemize}
\item The first argument has an index
\item The first argument is coindexed with \ft{spr} of the second argument
\item Both the first and  second arguments play 
semantic roles in the \prd{try} relation
\item Very little had to be stipulated in the entry 
\end{itemize}
\end{minipage}
\end{tabular}

\qaslide{Questions}{
\begin{itemize}
\item What rules out dummies and idiom chunks as 
subjects of \lex{try}?
  \answer{\\ {\small\emp{\lex{try}'s subject is NP$_i$, which fills the \textit{trier} role; dummies and idiom chunks have no index.}}}
\item What accounts for the semantic non-equivalence of 
pairs like the following?
\begin{xlisti}
  \ex \eng{Reporters tried to interview the candidate.}
  \ex \eng{The candidate tried to be interviewed by reporters.}
\end{xlisti}
  \answer{{\small\emp{The \ft{trier} role is filled by a different index in each: the reporters in (i), the candidate in (ii).  The \textbf{interview} predication is the same in both.}}}
\item Why does \lex{continue} behave differently in these 
respects?
  \answer{\\ {\small\emp{It is a raising verb: its subject is just its complement's \ft{spr}, with no role of its own, so whatever the complement allows is fine, and active and passive mean the same.}}}
\end{itemize}
}

\myslide{\lex{try} with an Active Complement}
\begin{minipage}[c]{0.58\textwidth}
\scalebox{0.72}{\begin{forest}
cat
[{S\avmbox{A}} [{\avmbox{1} NP$_i$} [{The police}, roof]]
[{VP \avm{[ \ft{spr} < \avmbox{1}$_i$ > ]}} [{V \avm{[ \ft{spr} < \avmbox{1}$_i$ > ]}} [{tried}]]
[{VP \avm{[ \ft{spr} < \avmbox{2}$_i$ > ]}} [{V \avm{[ \ft{spr} < \avmbox{2}$_i$ > ]}} [{to}]]
[{VP \avm{[ \ft{spr} < \avmbox{2}$_i$ > ]}} [{V \avm{[ \ft{spr} < \red{\avmbox{2}$_i$} > ]}} [{arrest}]]
[{NP$_j$} [{the suspects}, roof]]]]]]
\end{forest}}
\end{minipage}\hfill
\begin{minipage}[c]{0.4\textwidth}
\scalebox{0.8}{\avm{[ restr & \avmbox{A} <
  [ reln & \textbf{try} \\ sit  & s$_1$ \\ trier & \OV{i} \\ arg & s$_2$ ], \\
  [ reln & \textbf{arrest} \\ sit  & s$_2$ \\ arrester & \OV{i} \\ arrested & j],
  \ldots > ]}}
\end{minipage}

\emp{Two different NPs, \avmbox{1} and \avmbox{2}, sharing one index $i$:
  the police are the \ft{trier} and the \ft{arrester}.}

\myslide{\lex{try} with Passive Complement}
\begin{minipage}[c]{0.58\textwidth}
\scalebox{0.6}{\begin{forest}
cat
[{S\avmbox{A}} [{\avmbox{1} NP$_j$} [{The suspects}, roof]]
[{VP \avm{[ \ft{spr} < \avmbox{1}$_j$ > ]}} [{V \avm{[ \ft{spr} < \avmbox{1}$_j$ > ]}} [{tried}]]
[{VP \avm{[ \ft{spr} < \avmbox{2}$_j$ > ]}} [{V \avm{[ \ft{spr} < \avmbox{2}$_j$ > ]}} [{to}]]
[{VP \avm{[ \ft{spr} < \avmbox{2}$_j$ > ]}} [{V \avm{[ \ft{spr} < \avmbox{2}$_j$ > ]}} [{be}]]
[{VP \avm{[ \ft{spr} < \avmbox{2}$_j$ > ]}} [{V \avm{[ \ft{spr} < \red{\avmbox{2}$_j$} > ]}} [{arrested}]]
[{PP$_i$} [{P$_i$} [{by}]]
[{NP$_i$} [{the police}, roof]]]]]]]]
\end{forest}}
\end{minipage}\hfill
\begin{minipage}[c]{0.4\textwidth}
\scalebox{0.8}{\avm{[ restr & \avmbox{A} <
  [ reln & \textbf{try} \\ sit  & s$_1$ \\ trier & \OV{j} \\ arg & s$_2$ ], \\
  [ reln & \textbf{arrest} \\ sit  & s$_2$ \\ arrester & i \\ arrested & \OV{j}],
  \ldots > ]}}
\end{minipage}

\emp{Now the \ft{trier} is $j$, the suspects: a different
  \textbf{try} predication, so a different meaning from the active.}

\qaslide{ARG-ST of Raising vs.\ Control Verbs}{
\bigskip
\begin{center}
\begin{tabular}{cc}
  \avm{< NP$_i$,\ \ 
    VP [ inf & $+$ \\ spr  & < NP$_i$ > \\ index & s$_2$ ]
    >
}& 
  \avm{< \1 NP,\ \ 
    VP [ inf & $+$ \\ spr  & < \1 > \\ index & s$_2$ ]
    >
} \\
  \\[2ex]
  \answer{\txx{Control}} & \answer{\txx{Raising}}
\end{tabular}
\end{center}
\bigskip
\begin{itemize}
\item Which is which?  Why?
  \answer{\\ {\small\emp{Left: control.  Two NPs that merely share an index $i$, so the subject must have one (no dummies), and it can fill a role of the verb's own.  Right: raising.  One NP, \avmbox{1}, that is both the verb's subject and its complement's: whatever the complement wants, dummies included, is passed straight through.}}}
\end{itemize}
}

\section{Raising and Control in Transformational Grammar}

\myslide{Raising \& Control in Transformational Grammar}


\begin{itemize}
\item Raising
  \begin{exe}
    \ex    \eng{\rnode{A}{\rule{4em}{1pt}} continue [\rnode{B}{\ul{the dogs}} to bark]}
    
     \ncbar[offsetB=4pt,angleA=-90,nodesep=3pt]{<-}{A}{B}
  \end{exe}
\item Control
  \begin{exe}
    \ex \eng{[the dogs]$_i$ try [NP$_i$  to bark]}
  \end{exe}
  \begin{itemize}
  \item In early TG, the NP got deleted.
  \item In more recent TG, it’s a silent pronoun.
  \end{itemize}
\end{itemize}

\qaslide{Problems with the TG Accounts}{


\begin{itemize}
\item Details never fully worked out (e.g.,\ where does \lex{to}
come from?)

\item What blocks 
  \begin{xlisti}
    \ex \bad{The cat continued (for) the dog to bark}
    \ex \bad{The cat tried (for) the dog to bark}?

  \end{xlisti}
  \answer{{\small\emp{Nothing natural: the TG account needs an extra stipulation.  For us the complement is a VP whose \ft{spr} is the verb's subject, so there is no place for another subject.}}}

\item Failure of experimental attempts to find evidence for 
psychological reality of these transformations.
\end{itemize}
}

\section{Object Raising and Object Control}

\myslide{We Make Another Raising/Control Distinction}

\begin{exe}
  \ex 
  \begin{xlist}
    \ex \eng{I \ul{expected} Leslie to be aggressive.} \hfill \val{orv}
    \ex \eng{I \ul{persuaded} Leslie to be aggressive.} \hfill \val{ocv}
  \end{xlist}
\end{exe}

\begin{center} \small
  \begin{tabular}{cc}
\txx{Object-Raising Verb Lexeme} (\val{orv-lxm}) &
\txx{Object-Control Verb Lexeme} (\val{ocv-lxm}) \\
  \avm{[ arg-st & < NP, \1,
                   [ spr   & < \1 > \\
                      comps & < ~ > \\
                      index & \OV{s}  ] > \\
      sem & [ restr & < [ arg & s ] > ] ]
}
&    \avm{[ arg-st & < NP, NP$_i$,
                     [  spr & < NP$_i$ > \\
                        comps & < ~ > \\
                        index & \OV{s}  ] > \\
       sem & [ restr & < [ arg & s ] > ] ]
}
\end{tabular}
\end{center}
\begin{itemize}
\item The formal distinction is again between tagging and coindexing
\item This time it’s the second argument and the \ft{spr} of the third
  argument.
\end{itemize}

\myslide{Example \typ{orv-lxm} and \typ{ocv-lxm} Entries}

\begin{center} \small
  \begin{tabular}{cc}
\txx{Object-Raising Verb Lexeme} (\val{orv-lxm}) &
\txx{Object-Control Verb Lexeme} (\val{ocv-lxm}) \\
  \avm{\tiny
     < \textnormal{expect}, \ 
     [\val{\tiny orv-lxm} \\
     arg-st & < NP$_j$, X, VP [ inf & $+$ ] > \\
      sem & [ index & s \\
               restr & < [ reln & \prd{expect} \\
                             sit & s \\
                             expecter & j  ] > ] ] >
}
& \avm{\tiny
     < \textnormal{persuade}, \ 
     [\val{\tiny ocv-lxm} \\
     arg-st & < NP$_j$, NP$_i$, VP [ inf & $+$ ] > \\
      sem & [ index & s \\
               restr & < [ reln & \prd{persuade} \\
                             sit & s \\
                             persuader & j\\
                             persuadee & i ] > ] ] >
} 
\end{tabular}
\end{center}



\begin{itemize}
\item Note that the \prd{persuade} relation has three arguments,
  but the \prd{expect} relation has only two
\item And the object’s \ft{index} plays a role in the \prd{persuade} relation,
  but not in the \prd{expect} relation
\end{itemize}

\myslide{Aside: In a Big Grammar}
\MyLogo{English Resource Grammar (ERG), lexical types for \lex{to}-VP complements}
\begin{itemize}\addtolength{\itemsep}{-1ex}
\item The same two ingredients, a shared tag or a shared index, give
  dozens of lexical types in a broad-coverage grammar of English:
\end{itemize}
\begin{center}\small
\begin{tabular}{lll}
  subject raising & \iz{v\_vp\_ssr\_le} & \eng{There failed to be a link.} \\
  subject control & \iz{v\_vp\_seq\_le} & \eng{Kim intended to win.} \\
  object raising & \iz{v\_np-vp\_sor\_le} & \eng{Kim expected Sandy to win.} \\
  object control & \iz{v\_np-vp\_oeq\_le} & \eng{Kim invited Sandy to stay.} \\
  subject control, with object & \iz{v\_np-vp\_aeq\_le} & \eng{Kim promised Sandy to stay.} \\
  raising adjective & \iz{aj\_vp\_i-ssr\_le} & \eng{There are bound to be problems.} \\
  control adjective & \iz{aj\_pp-vp\_i-it\_le} & \eng{It is easy for Kim to win.} \\
  control noun & \iz{n\_vp\_c\_le} & \eng{Kim has the ability to win.} \\
\end{tabular}
\end{center}
\begin{itemize}\addtolength{\itemsep}{-1ex}
\item \iz{ssr} = subject-to-subject raising, \iz{sor} = subject-to-object raising,
  \iz{seq}/\iz{oeq} = subject/object ``equi'' (the old name for control).
\item Problem Set 10 asks you to classify a few of these yourself.
\end{itemize}

% \myslide{Justifying the difference between expect and persuade (Prob. 12.4)}

% Construct examples of each of the following four 
% types which show a contrast between expect and 
% persuade:
% \begin{enumerate}
% \item Examples with dummy \lex{there}
% \item Examples with dummy \lex{it}
% \item Examples with idiom chunks
% \item Examples of relevant pairs of sentences containing 
% active and passive complements.  Indicate whether 
% they are or are not paraphrases of each other. 
% \end{enumerate}

\section{Raising and Control in Czech}

\myslide{Does Czech Have the Same Split?}
\MyLogo{You are the native speakers: check every judgement here}
\begin{itemize}
\item Czech also has NP-V-\textsc{inf} constructions with two kinds of verb:
  \begin{center}
  \begin{tabular}{lp{0.75\textwidth}}
    \txx{Raising?} & \lex{začít} `begin', \lex{přestat} `stop', \lex{zdát se}
    `seem', \lex{hrozit} `threaten', \lex{muset} `must' \\[1ex]
    \txx{Control?} & \lex{snažit se} `try', \lex{pokusit se} `attempt',
    \lex{chtít} `want', \lex{odmítnout} `refuse', \lex{slíbit} `promise' \\
  \end{tabular}
  \end{center}
\item Slovak is parallel: \lex{začať}, \lex{prestať}, \lex{zdať sa} vs.\
  \lex{snažiť sa}, \lex{pokúsiť sa}, \lex{sľúbiť}.
\end{itemize}

\qaslide{Which Tests Still Work?}{
\begin{itemize}
\item Czech has no dummy subject: no \lex{there}, no \lex{it}.
  \eng{Prší} `(it) rains' has no subject at all.
\item So what happens to our four tests?
  \answer{
  \begin{itemize}
  \item \emp{dummy \lex{there}, dummy \lex{it} \into\ one test with
      \txx{subjectless} impersonal verbs}
  \item \emp{idiom chunk subjects: works as in English}
  \item \emp{active/passive paraphrase: works as in English}
  \end{itemize}}
\end{itemize}
}

\myslide{Test 1: Impersonal Verbs with \lex{začít}, \lex{přestat}}
\begin{exe}
\ex \gll Začalo pršet. \\
    began.\gm{n.sg} rain.\gm{inf} \\
\trans `It began to rain.'
\ex \gll Přestalo sněžit. \\
    stopped.\gm{n.sg} snow.\gm{inf} \\
\trans `It stopped snowing.'
\ex \gll Začíná jít o život. \\
    begins go.\gm{inf} about life \\
\trans `It is becoming a matter of life and death.'
\end{exe}
\begin{itemize}
\item No subject anywhere, and the verb shows default neuter singular agreement.
\end{itemize}

\qaslide{Test 1: Impersonal Verbs with \lex{snažit se}, \lex{slíbit}}{
\begin{exe}
\ex \gll *Pokusilo se pršet. \\
    attempted.\gm{n.sg} \gm{refl} rain.\gm{inf} \\
\ex \gll *Snažilo se sněžit. \\
    tried.\gm{n.sg} \gm{refl} snow.\gm{inf} \\
\ex \gll *Slíbilo jít o život. \\
    promised.\gm{n.sg} go.\gm{inf} about life \\
\end{exe}
\begin{itemize}
\item Why the difference?
  \answer{\\ \emp{A subjectless complement means a subjectless matrix
      verb.  \lex{začít} passes the (empty) subject requirement
      through, so it is fine.  \lex{snažit se} needs a subject with an
      index to be the trier, and there is nothing to be it.}}
\end{itemize}
}

\qaslide{Test 2: Idiom Chunks (\eng{nervy})}{
\begin{exe}
\ex \gll Tečou mu nervy. \\
    flow him.\gm{dat} nerves \\
\trans `He is losing his patience.' (lit.\ `his nerves are flowing')
\ex \gll Začaly mu téct nervy. \\
    began.\gm{pl} him.\gm{dat} flow.\gm{inf} nerves \\
\trans `He started losing his patience.'
\ex \gll \#Nervy se mu snažily téct. \\
    nerves \gm{refl} him.\gm{dat} tried.\gm{pl} flow.\gm{inf} \\
\trans only literal: `The nerves tried to flow.'
\end{exe}
\begin{itemize}
\item Why does the idiom survive \lex{začít} but not \lex{snažit se}?
  \answer{\\ \emp{\lex{začít} gives its subject no role, so \eng{nervy}
      can stay a piece of the idiom.  \lex{snažit se} makes its subject
      the trier, and nerves cannot try.}}
\end{itemize}
}

\qaslide{Test 2: Idiom Chunks (\eng{síly})}{
\begin{exe}
\ex \gll Docházejí mu síly. \\
    run.out him.\gm{dat} strength.\gm{pl} \\
\trans `He is running out of strength.'
\ex \gll Začaly mu docházet síly. \\
    began.\gm{pl} him.\gm{dat} run.out.\gm{inf} strength.\gm{pl} \\
\trans `He started running out of strength.'
\ex \gll \#Síly mu slíbily docházet. \\
    strength.\gm{pl} him.\gm{dat} promised.\gm{pl} run.out.\gm{inf} \\
\trans only literal: `The strength promised to run out.'
\end{exe}
\begin{itemize}
\item Same pattern.  Which type is \lex{slíbit} `promise'?
  \answer{\\ \emp{Control: its subject is the promiser, so it needs an
      index, and an idiom chunk has none.}}
\end{itemize}
}

\myslide{Test 3: Passive Paraphrase with \lex{začít}}
\begin{exe}
\ex \gll Lékaři začali vyšetřovat pacienta. \\
    doctors began.\gm{pl} examine.\gm{inf} patient.\gm{acc} \\
\trans `The doctors began to examine the patient.'
\ex \gll Pacient začal být vyšetřován lékaři. \\
    patient began.\gm{sg} be.\gm{inf} examined doctors.\gm{ins} \\
\trans `The patient began to be examined by the doctors.'
\end{exe}
\begin{itemize}
\item Do these describe the same situations?
\end{itemize}

\qaslide{Test 3: Passive Paraphrase with \lex{snažit se}}{
\begin{exe}
\ex \gll Lékaři se snažili vyšetřit pacienta. \\
    doctors \gm{refl} tried.\gm{pl} examine.\gm{inf} patient.\gm{acc} \\
\trans `The doctors tried to examine the patient.'
\ex \gll Pacient se snažil být vyšetřen lékaři. \\
    patient \gm{refl} tried.\gm{sg} be.\gm{inf} examined doctors.\gm{ins} \\
\trans `The patient tried to be examined by the doctors.'
\end{exe}
\begin{itemize}
\item And these?
  \answer{\\ \emp{No.  With \lex{začít} both sentences say the examining
      began: beginning is a property of the situation, however it is
      phrased.  With \lex{snažit se} the trier changes from the doctors
      to the patient, so the meanings differ.}}
\end{itemize}
}

\qaslide{Czech Bonus: Agreement}{
\begin{exe}
\ex \gll Děti začaly spát. \\
    children began.\gm{pl} sleep.\gm{inf} \\
\trans `The children began to sleep.'
\end{exe}
\begin{itemize}
\item \lex{začít} agrees with a subject it gives no role to.  Why is
  this exactly what tagging predicts?
  \answer{\\ \emp{The raising verb's \ft{spr} \emph{is} its complement's
    \ft{spr}, tag \avmbox{1}, and the SHAC makes the verb agree with
    its \ft{spr}, role or no role.  With no subject at all, as in
    \eng{Začalo pršet}, the agreement falls back to neuter singular.}}
\end{itemize}
}

\qaslide{Czech Bonus: Objects}{
\begin{itemize}
\item Object control exists with infinitives, with accusative or dative
  objects:
  \begin{exe}
  \ex \gll Naučil jsem Petra plavat. \\
      taught \gm{aux.1sg} Petr.\gm{acc} swim.\gm{inf} \\
  \trans `I taught Petr to swim.'
  \ex \gll Dovolil jsem Petrovi odejít. \\
      allowed \gm{aux.1sg} Petr.\gm{dat} leave.\gm{inf} \\
  \trans `I allowed Petr to leave.'
  \end{exe}
\item But \lex{expect}-type object raising needs a finite clause:
  \eng{Očekával jsem, že Petr přijde}, not \bad{Očekával jsem Petra přijít}.
  What does that mean for \val{orv-lxm} in Czech?
  \answer{\\ \emp{Czech simply has no verb of this type with an
      infinitive complement.  A gap in the lexicon, not in the theory.}}
\end{itemize}
}

\section{Problem Set 10}

\myslide{P1: Classifying Verbs}
\MyLogo{Based on \citet[Chapter 12, Problem 1]{sag:wasow:bender:2003}}
\small

\noindent Classify the following verbs as raising or control:
%Give as many distinct tests as you can to justify your
%classifications, constructing three or four interestingly different
%examples for each verb.

\begin{itemize}
\item \lex{tend, decide, manage, fail, happen}
\end{itemize}

\noindent
Justify your classification by applying each of the following four
tests to each verb.  Show your work by providing relevant examples
and indicating their grammaticality. 

\begin{itemize}
\item[(i)] Can the verb take a  dummy \lex{there} subject if and only if its complement selects for a dummy \lex{there} subject?

\item[(ii)] Can the verb take a dummy \lex{it} subject if and only if
  its complement selects for a dummy \lex{it} subject?

\item[(iii)] Can the verb take an idiom chunk subject if and only
if the rest of the idiom is in its complement?

\item[(iv)] Do pairs of sentences containing active and 
passive complements to the verb end up being  paraphrases 
of each other?
\end{itemize}

\noindent
Make sure to restrict your attention to cases of the form: \hbox{NP V
  {\it to} VP}. That is, ignore cases like \eng{Kim manages a store},
\eng{Alex failed physics}, and any other valence that doesn't resemble
the \lex{continue} vs.\ \lex{try} pattern.
\normalsize

\myslide{P2: Classifying \index{adjective}Adjectives}
\MyLogo{Based on \citet[Chapter 12, Problem 2]{sag:wasow:bender:2003}}
Classify the following adjectives as \index{raising!adjective}raising or 
\index{control!adjective}control:
%Justify your classifications, constructing three or four
%interestingly different relevant examples for each adjective.

\begin{itemize}
\item \lex{anxious, apt, bound, certain, lucky}
\end{itemize}

\noindent
Justify your classification by providing each of the four types
of data discussed in the previous problem for each adjective.

Make sure to restrict your attention to cases of the form: \hbox{NP
  \lex{be} Adj \lex{to} VP}. That is, ignore cases like \eng{Kim is
  anxious about the exam}, \eng{Carrie is certain of the answer}, and
any other valence that doesn't resemble the \lex{likely} vs.\
\lex{eager} pattern.

\myslide{P3: \lex{expect} vs.\ \lex{persuade}}
\MyLogo{Based on \citet[Chapter 12, Problem 4]{sag:wasow:bender:2003}}

Construct the arguments that underlie the proposed distinction between
\typ{orv-lxm} and \typ{ocv-lxm}.

Construct examples of each of the following four types that 
show a contrast between \lex{expect} and \lex{persuade}.  Explain how
the contrasts are accounted for by the differences in the types
\typ{orv-lxm} and \typ{ocv-lxm} and/or the lexical entries for \lex{expect}
and \lex{persuade}.

\begin{itemize} \addtolength{\itemsep}{-1ex}
\item[(i)] Examples with dummy \lex{there}.

\item[(ii)] Examples with dummy \lex{it}.

\item[(iii)] Examples with idiom chunks.

\item[(iv)] Examples of relevant pairs of sentences containing 
active and passive complements.  Indicate whether they are or are
not paraphrases of each other.
\end{itemize}



% \noindent Again, make sure you ignore all irrelevant uses of
% these verbs, including cases of 
% %finite 
% CP complements, e.g.\ \eng{persuade NP that ...} or
% \eng{expect that ...} and anything else not directly relevant (\eng{I
%   expect to go}, \eng{I am expecting Kim}, \eng{She is expecting}, and
% so forth).

\myslide{P4: A Type for Existential \lex{be}}
\MyLogo{Based on \citet[Chapter 12, Problem 6]{sag:wasow:bender:2003}}

The \lex{be} that takes \lex{there} as its subject wasn't given a true
lexical type in Chapter 11 \citep{sag:wasow:bender:2003}, because no
suitable type had been introduced.  One of the types in this chapter
will do, if we make some of its constraints defeasible.

\begin{itemize}
\item[A.] Which of the types introduced in this chapter comes
closest to being consistent with the constraints on \lex{there}-taking
\lex{be}?
\item[B.] Rewrite that type indicating which constraints must
be made defeasible.
\item[C.] Give a streamlined lexical entry for the \lex{there}-taking
\lex{be} that stipulates only those constraints that are truly
idiosyncratic to the lexeme.
\end{itemize}




\myslide{Overview}
\MyLogo{\citet{sag:wasow:bender:2003} --- Chapter 12}
\begin{itemize}
\item Introduction: two kinds of NP-V-(NP)-\lex{to}-VP
\item Infinitival \lex{to}
\item Subject-raising verbs
\item Subject-control verbs
\item Raising and control in Transformational Grammar
\item Object raising and object control
\item Raising and control in Czech
\item Problem Set 10: SWB 12.1, 12.2, 12.4, 12.6
\end{itemize}



%\input{lec-10-control-q}


\myslide{Acknowledgments and References}

\begin{itemize}
\item Course design and slides borrow heavily from Emily Bender's course:
\textit{Linguistics 566: Introduction to Syntax for Computational Linguistics}
\\ \url{http://courses.washington.edu/ling566}
\item This lecture complements \citet[Chapter 12]{sag:wasow:bender:2003}
\end{itemize}
\bibliographystyle{aclnat}
\bibliography{abb,mtg,nlp,ling}
\end{document}

%%% Local Variables: 
%%% coding: utf-8
%%% mode: latex
%%% TeX-PDF-mode: t
%%% TeX-engine: xetex
%%% End: 
